% chapter4_system_design.tex
% BondHealth Healthcare Platform - System Design Chapter

\chapter{System Design}

\justifying
This chapter presents the comprehensive system design of the BondHealth healthcare platform,
detailing the architectural decisions, data flow, database schema, API design, and module
specifications. The system design follows industry best practices for healthcare applications,
ensuring scalability, security, maintainability, and seamless connectivity between patients,
doctors, hospital administrators, and laboratory technicians.

% =============================================================================
% 4.1 SYSTEM ARCHITECTURE OVERVIEW
% =============================================================================

\section{System Architecture Overview}

\justifying
The BondHealth platform is built on a robust three-tier architecture ensuring scalability,
maintainability, and security while providing seamless connectivity between all stakeholders,
with clear separation of concerns between presentation, business logic, and data persistence layers.
Each tier is designed to operate independently while communicating through well-defined interfaces,
allowing any single layer to be scaled or modified without disrupting the others. The three tiers
are the Presentation Layer, the Application Layer, and the Data Layer, each of which is described in
detail below.

\begin{figure}[ht]
    \centering
    \includegraphics[width=0.9\textwidth]{images/chapter4/architecture.png}
    \caption{Three-Tier System Architecture of BondHealth Platform}
    \label{fig:ch4-architecture}
\end{figure}

\subsection{Presentation Layer (Client Tier)}

\justifying
The presentation layer serves as the user interface through which all stakeholders interact
with the platform. It utilizes server-side rendered HTML with Tailwind CSS for responsive,
mobile-first design. Key advantages include fast initial page loads, improved SEO, enhanced
security by keeping business logic on the server, and a consistent experience across all
devices and screen sizes.

\noindent
\justifying
The interface is organized into role-specific dashboards: patient dashboards focus on
appointment management and medical record access, doctor dashboards prioritize clinical
workflow efficiency, admin dashboards provide comprehensive hospital management, and lab
technician dashboards streamline test processing and result reporting.

\subsection{Application Layer (Business Tier)}

\justifying
The application layer implements all business logic, validation, workflow orchestration, and
security controls using Node.js with the Express framework. It exposes RESTful API endpoints
consumed by the presentation layer and enforces security at every request boundary. JWT-based
authentication middleware validates tokens on each protected route, while Role-Based Access
Control (RBAC) ensures that users can only perform operations permitted by their assigned role.
A dedicated request validation layer rejects malformed or incomplete payloads before they reach
any business logic, and centralized error handling ensures that all failure modes return consistent,
structured responses. Request rate limiting is applied at the router level to protect against
denial-of-service conditions.

\subsection{Data Layer (Database Tier)}

\justifying
The data layer manages all persistent storage using Neon PostgreSQL, a serverless solution
offering automatic scaling and high availability. Connections are managed through the
\texttt{@neondatabase/serverless} pooling adapter, which eliminates per-request connection
overhead while respecting Neon's serverless concurrency limits. All primary keys use UUID v4
to support distributed record generation without coordination. Native PostgreSQL array types
are used for multi-value attributes, foreign key constraints enforce referential integrity
across all relationships, and SSL encryption is required for every database connection.

% =============================================================================
% 4.2 DATA FLOW DIAGRAMS
% =============================================================================

\section{Data Flow Diagrams}

\justifying
Data Flow Diagrams (DFDs) provide a graphical representation of how information moves
through the BondHealth platform. They capture the inputs, transformations, and outputs at each
stage of a workflow, making it possible to reason about correctness, latency, and failure
modes before implementation. This section describes the three most critical workflows in the
system: appointment booking, lab report processing, and payment processing.

\subsection{Appointment Booking Data Flow}

\justifying
The appointment booking process enables patients to schedule consultations with doctors based
on real-time availability. The workflow begins when the patient authenticates via JWT and
proceeds to search and filter doctors by specialization and availability. Once a suitable
doctor is identified, the patient selects a time slot; database transactions prevent
double-booking by locking the slot for the duration of the operation. The patient then
completes payment through the integrated gateway, after which confirmation is dispatched via
email, SMS, and in-app notifications. Both the patient and doctor dashboards are updated
immediately upon booking, ensuring that neither party sees stale schedule information.

\begin{figure}[ht]
    \centering
    \includegraphics[width=0.95\textwidth]{images/chapter4/appointment-flow.png}
    \caption{Appointment Booking Data Flow Diagram}
    \label{fig:chapter4-appointment-flow}
\end{figure}

\subsection{Lab Report Processing Data Flow}

\justifying
Lab technicians authenticate with a \texttt{lab} role and locate patients via UUID, name, or
QR code. They then select test types from a predefined list and enter results, which are
validated in real time against reference ranges; critical values trigger immediate alerts to
the responsible clinician. A priority level is assigned to each report, and any supporting
files are virus-scanned before being written to storage. Sharing preferences are configured
at submission time, and once the report is submitted the system dispatches notifications
scaled to the assigned priority level, ensuring urgent results reach the patient and doctor
without delay.

\subsection{Payment Processing Data Flow}

\justifying
The payment workflow is triggered upon appointment booking or post-consultation billing. The
system calculates the total amount due and presents the patient with a choice of payment
methods, including card, insurance, digital wallet, and UPI. The transaction is routed to
the appropriate external gateway using 3D Secure authentication to satisfy card-scheme
requirements, and each transaction is recorded with a unique UUID alongside the gateway's
own reference identifier. On a successful response, the appointment status is updated to
confirmed, an invoice PDF is generated, and a receipt is emailed to the patient.
Asynchronous gateway events, such as delayed bank confirmations, are handled through
dedicated webhook endpoints that reconcile the transaction state without requiring the user
to remain on the page.

% =============================================================================
% 4.3 DATABASE SCHEMA DESIGN
% =============================================================================

\section{Database Schema Design (Neon PostgreSQL)}

\justifying
The BondHealth platform utilizes Neon PostgreSQL, providing automatic scaling, high
availability, and connection pooling. The schema is the authoritative definition of all data
the system stores and the relationships between those data. Every design decision in the
schema — from key types to constraint styles to timestamp conventions — has been made to
support the platform's clinical workflows reliably while remaining maintainable as the
system evolves. The schema uses UUID primary keys, array data types, and comprehensive
foreign key constraints.

% CHANGE 2: ER Diagram figure placed first
\begin{figure}[ht]
    \centering
    \includegraphics[width=\textwidth]{images/chapter4/schema.png}
    \caption{Database Entity-Relationship Diagram}
    \label{fig:schema}
\end{figure}

\justifying
\noindent
The Entity-Relationship diagram above illustrates the complete relational structure of
the BondHealth database. At the center is the \texttt{users} table, which serves as the
authentication anchor for all four system roles. Each role-specific table references
\texttt{users} via a foreign key on \texttt{user\_id}, providing a unified identity per
account and eliminating separate credential stores per role.

\justifying
\noindent
The \texttt{hospitals} table acts as the second major anchor, grouping doctors, medicines,
and lab technicians under a specific facility. The \texttt{appointments} table links
patients and doctors with scheduling, status, and payment information. The
\texttt{lab\_reports} and \texttt{prescriptions} tables similarly bridge patients and
doctors, capturing clinical outputs. The \texttt{doctor\_leave} table integrates with
appointment availability logic, the \texttt{medicines} table is scoped per hospital for
independent inventory management, and the \texttt{payments} table maintains a one-to-one
relationship with appointments covering the full financial transaction lifecycle.

\justifying
\noindent
The schema follows third normal form (3NF), uses \texttt{created\_at} and
\texttt{updated\_at} timestamps on all tables for audit compliance, applies strategic
indexes on frequently queried columns, and uses UUID primary keys with serverless
PostgreSQL for horizontal scalability. Sensitive fields such as \texttt{password\_hash}
are never exposed through API responses.

\justifying
\noindent
Referential integrity is enforced throughout the schema using foreign key constraints
with carefully chosen cascade rules. Deleting a \texttt{users} record cascades to the
corresponding role-specific profile row, ensuring no orphaned patient or doctor records
remain after account removal. By contrast, \texttt{appointments} and \texttt{lab\_reports}
use a restrict rule on their patient and doctor foreign keys, preventing accidental data
loss for clinically significant records. This combination of cascade and restrict rules
reflects the differing retention requirements across the schema.

\justifying
\noindent
PostgreSQL's native array type is used for multi-value attributes such as a doctor's
available weekdays and a patient's known allergies or chronic conditions. Storing these
as arrays rather than in separate junction tables reduces join complexity for the most
frequent read queries, while still supporting array-specific operators such as
\texttt{@>} (contains) for availability filtering.

\justifying
\noindent
Status fields across the schema use constrained string enumerations rather than integer
codes, making query results self-documenting and reducing the risk of application-level
mapping errors. For example, \texttt{appointments.status} accepts only the values
\texttt{pending\_payment}, \texttt{confirmed}, \texttt{completed}, and \texttt{cancelled},
enforced through a \texttt{CHECK} constraint. Similarly, \texttt{doctor\_leave.status}
is constrained to \texttt{pending}, \texttt{approved}, and \texttt{rejected}, ensuring
that the leave approval workflow cannot produce an unrecognised state in the database.

\justifying
\noindent
All timestamp columns use \texttt{TIMESTAMPTZ} rather than a timezone-naive type,
storing values in UTC and allowing the application layer to convert to the user's local
timezone at render time. This is particularly important for a platform intended to support
multiple hospital facilities, where doctors and patients may operate across different
time zones. Consistent UTC storage eliminates ambiguity in scheduling queries and ensures
that appointment slot comparisons remain correct regardless of the server's local
timezone configuration.

\justifying
\noindent
The schema is managed through a versioned migration system, capturing every structural
change as an incremental script across development, staging, and production environments.
New hospitals are provisioned without structural changes --- seed scripts initialise default
roles, department categories, and configuration records for each facility.

\justifying
\noindent
The schema comprises ten core tables, each scoped to a distinct domain of the platform.
The \texttt{users} and \texttt{hospitals} tables form the two anchoring entities from which
all other tables derive their identity and organisational context. Clinical output tables
such as \texttt{lab\_reports} and \texttt{prescriptions} depend on both a patient and a
doctor reference, reflecting the bilateral nature of medical record ownership. Operational
tables such as \texttt{doctor\_leave}, \texttt{medicines}, and \texttt{payments} each serve
a narrower purpose but are equally critical to the platform's day-to-day workflows.
Together, these ten tables cover every data entity required by the four user roles without
redundancy, and their relationships have been designed to make the most common read queries
--- patient history, doctor availability, appointment lookup --- executable with a minimal
number of joins.

\justifying
\noindent
Each table follows a consistent internal convention. Every row is identified by a UUID
primary key generated at the database level, every table carries \texttt{created\_at} and
\texttt{updated\_at} timestamp columns for audit traceability, and every foreign key
relationship is accompanied by an explicit constraint with a deliberate cascade or restrict
rule. Status and role fields use \texttt{CHECK}-constrained string values rather than
numeric codes, so that the contents of any row are self-explanatory when inspected directly
in the database. 

\noindent
The following table provides a concise summary of each core table, its purpose, and
its primary fields as implemented in the Neon PostgreSQL schema.

\begin{table}[H]
    \small
    \centering
    \renewcommand{\arraystretch}{1.3}
    \begin{tabular}{|>{\raggedright\arraybackslash}p{2.6cm}|>{\raggedright\arraybackslash}p{6.3cm}|>{\raggedright\arraybackslash}p{4.7cm}|}
        \hline
        \textbf{Table Name} & \textbf{Description} & \textbf{Primary Fields} \\
        \hline
        users
            & Authentication credentials for all system users.
            & user\_id (UUID), username, email, password\_hash, role \\
        \hline
        hospitals
            & Healthcare facility information.
            & hospital\_id (UUID), name, type, city, phone \\
        \hline
        patients
            & Patient medical info and emergency contacts.
            & patient\_id (UUID), user\_id, full\_name, blood\_type, conditions[] \\
        \hline
        doctors
            & Professional details and availability.
            & doctor\_id (UUID), user\_id, hospital\_id, specialization, fee \\
        \hline
        appointments
            & Patient-doctor appointment records.
            & appointment\_id (UUID), patient\_id, doctor\_id, date, time, status \\
        \hline
        lab\_reports
            & Laboratory test results with file attachments.
            & report\_id (UUID), patient\_id, doctor\_id, test\_type, results \\
        \hline
        prescriptions
            & Medication prescriptions and refill info.
            & prescription\_id (UUID), patient\_id, doctor\_id, medicine\_name \\
        \hline
        doctor\_leave
            & Doctor leave applications and approvals.
            & leave\_id (UUID), doctor\_id, from\_date, to\_date, reason \\
        \hline
        medicines
            & Pharmacy inventory with stock levels.
            & medicine\_id (UUID), hospital\_id, name, quantity, price \\
        \hline
        payments
            & Financial transactions and gateway references.
            & payment\_id (UUID), appointment\_id, patient\_id, amount, status \\
        \hline
    \end{tabular}
    \caption{Core Database Tables Schema}
    \label{tab:database-schema}
\end{table}

\vspace{0.3cm}
\noindent
\textbf{Database Indexes for Performance Optimization.}
Strategic indexes are applied to the columns most frequently used in \texttt{WHERE} clauses,
join conditions, and sort operations. \texttt{CREATE INDEX idx\_users\_email ON users(email)}
and \texttt{CREATE INDEX idx\_users\_username ON users(username)} accelerate login and user
lookup operations. \texttt{CREATE INDEX idx\_patients\_user\_id ON patients(user\_id)} enables
fast patient profile retrieval from a user identifier, while
\texttt{CREATE INDEX idx\_doctors\_hospital\_id ON doctors(hospital\_id)} and
\texttt{CREATE INDEX idx\_doctors\_specialization ON doctors(specialization)} support rapid
doctor discovery by facility and clinical specialty respectively. On the appointments side,
\texttt{CREATE INDEX idx\_appointments\_patient\_id ON appointments(patient\_id)},
\texttt{CREATE INDEX idx\_appointments\_doctor\_id ON appointments(doctor\_id)}, and
\texttt{CREATE INDEX idx\_appointments\_date ON appointments(appointment\_date)} together
ensure that retrieving a patient's history, a doctor's schedule, and date-range queries
all execute without full table scans.

% =============================================================================
% 4.4 API DESIGN
% =============================================================================

\section{API Design (RESTful Endpoints)}

\justifying
\noindent
The BondHealth platform implements a comprehensive RESTful API following industry best
practices, serving as the communication backbone between the presentation layer and business
logic. Endpoints are resource-oriented, use plural nouns for collections, appropriate HTTP
methods, and consistent JSON response formats. All protected endpoints require a valid JWT
token in an HTTP-only cookie, preventing exposure to client-side JavaScript.

\noindent
The API is organised into logical resource groups mirroring core domain entities: authentication,
patient management, doctor operations, appointment scheduling, laboratory workflows, prescription
handling, and administration. Middleware for role verification, request validation, and rate
limiting is scoped precisely at the router level. Incoming payloads are validated for type
correctness, required fields, and domain constraints before any business logic executes.

\noindent
All list-returning endpoints support \texttt{page} and \texttt{limit} query parameters,
returning a standardised envelope with the data array, total record count, and current page
metadata. Caching headers are applied selectively to read-heavy endpoints to reduce redundant
database queries without sacrificing freshness for time-sensitive resources such as appointment
slots. API versioning through \texttt{/api/v1/\ldots} supports future changes without breaking
existing clients.

\noindent
Error responses follow a consistent JSON structure containing an error code, a human-readable
message, and optional field-level validation details, enabling the frontend to surface
meaningful feedback to users without additional parsing logic. All endpoints return appropriate
HTTP status codes: \texttt{200} for successful retrieval, \texttt{201} for resource creation,
\texttt{400} for validation failures, \texttt{401} for unauthenticated requests, \texttt{403}
for authorisation violations, and \texttt{404} when a requested resource does not exist.
This uniformity means that the frontend can apply a single error-handling interceptor across
all API calls rather than implementing bespoke handling per endpoint, keeping the client-side
codebase lean and reducing the surface area for inconsistent user-facing error messages.

\noindent
The API also enforces per-user and per-IP rate limiting on sensitive endpoints such as login,
registration, and payment initiation, reducing exposure to brute-force and denial-of-service
attacks. For endpoints that trigger background operations --- such as appointment confirmation
and lab report notification --- responses are returned immediately with a \texttt{202 Accepted}
status while the background task completes asynchronously, keeping perceived response times
low under load. This distinction between synchronous and asynchronous response patterns is
applied deliberately: only operations whose background tasks are non-critical to the immediate
user action are deferred, while operations such as appointment slot locking and payment intent
creation remain fully synchronous to guarantee that the user receives a definitive outcome
before leaving the page. Taken together, the consistent status code semantics, structured
error envelopes, and mixed synchrony model give the API a predictable contract that both
the server-rendered EJS templates and any future client applications can depend on without
requiring changes to the underlying endpoint logic. The table below lists the core endpoints
along with their access requirements.

\begin{table}[H]
    \small
    \centering
    \renewcommand{\arraystretch}{1.3}
    \begin{tabular}{|>{\raggedright\arraybackslash}p{1.6cm}|>{\raggedright\arraybackslash}p{5.1cm}|>{\raggedright\arraybackslash}p{4.2cm}|>{\raggedright\arraybackslash}p{2.9cm}|}
        \hline
        \textbf{Method} & \textbf{Endpoint} & \textbf{Description} & \textbf{Access} \\
        \hline
        POST & /api/auth/register & New user registration & Public \\
        \hline
        POST & /api/auth/login & User authentication & Public \\
        \hline
        POST & /api/auth/logout & End user session & Auth \\
        \hline
        GET  & /api/auth/me & Current user profile & Auth \\
        \hline
        GET  & /api/patient & Patient profile data & Patient \\
        \hline
        GET  & /api/appointments & List user appointments & Patient/\allowbreak{}Doctor \\
        \hline
        POST & /api/appointments & Create new appointment & Patient \\
        \hline
        PUT  & /api/appointments/\allowbreak:id/\allowbreak reschedule & Change appointment time & Patient/\allowbreak{}Doctor \\
        \hline
        PUT  & /api/appointments/\allowbreak:id/\allowbreak cancel & Cancel appointment & Patient/\allowbreak{}Doctor \\
        \hline
        GET  & /api/doctors & Search available doctors & Patient \\
        \hline
        GET  & /api/doctors/\allowbreak:id/\allowbreak availability & Doctor's free slots & Patient \\
        \hline
        GET  & /api/doctor & Doctor profile & Doctor \\
        \hline
        GET  & /api/appointments/today & Today's schedule & Doctor \\
        \hline
        GET  & /api/lab-reports & Pending lab reports & Doctor \\
        \hline
        POST & /api/lab-reports & Create lab report & Lab \\
        \hline
        GET  & /api/hospital/data & Hospital statistics & Admin \\
        \hline
        POST & /api/hospital/\allowbreak add/doctor & Register new doctor & Admin \\
        \hline
        POST & /api/hospital/\allowbreak add/medicine & Add to inventory & Admin \\
        \hline
        POST & /api/payments/\allowbreak create-intent & Payment initialization & Patient \\
        \hline
        GET  & /api/messages/\allowbreak conversations & Chat history & Patient/\allowbreak{}Doctor \\
        \hline
        POST & /api/messages & Send message & Patient/\allowbreak{}Doctor \\
        \hline
    \end{tabular}
    \caption{Core API Endpoints}
    \label{tab:api-endpoints}
\end{table}

% =============================================================================
% 4.5 MODULE-WISE DETAILED DESIGN
% =============================================================================

\section{Module-wise Detailed Design}

\justifying
This section provides design specifications for each major module of the BondHealth platform,
describing their component structure and interactions with other modules. Each module encapsulates
a coherent set of responsibilities and exposes a well-defined interface to the rest of the system,
making it possible to develop, test, and evolve individual modules independently. The six modules
covered here are Authentication, the Patient Dashboard, the Doctor Dashboard, the Admin Dashboard,
the Lab Technician interface, and Appointment Booking.

\subsection{Authentication Module}

\justifying
The authentication module is the security gateway for the platform, responsible for verifying
user identities, managing sessions, and enforcing access controls. At its core, a JWT Token
Service generates tokens signed with HS256 that expire after seven days; these tokens are
delivered exclusively through HTTP-only cookies carrying \texttt{HttpOnly}, \texttt{Secure},
and \texttt{SameSite=Strict} flags, making them inaccessible to client-side JavaScript and
resistant to cross-site request forgery. Server-side session state is maintained in a Redis
store, enabling instant revocation without waiting for token expiry. Passwords are hashed
with bcrypt using ten salt rounds before storage, and RBAC middleware enforces the principle
of least privilege on every protected route by comparing the authenticated user's role against
the permissions required by the requested resource.

\subsection{Patient Dashboard Module}

\justifying
The patient dashboard provides a unified interface for managing all aspects of a patient's
healthcare journey. It dynamically loads only data relevant to the logged-in patient and
highlights pending actions such as unpaid bills, upcoming appointments, and new lab reports,
ensuring that time-sensitive items are immediately visible without requiring the patient to
navigate across multiple pages. The dashboard is composed of a Profile Summary Card, an
Upcoming Appointments section, a Past Appointments section, a tabbed Medical Records panel,
a Prescription Management view, Quick Action Buttons for the most common tasks, and a
Notification Center that distinguishes between read and unread items to help patients manage
their follow-up actions efficiently.

\subsection{Doctor Dashboard Module}

\justifying
The doctor dashboard is designed for clinical workflow efficiency, surfacing the information
and controls a clinician needs within a single view. A Dashboard Header displays quick
statistics and a status toggle that allows the doctor to mark themselves as available or
unavailable. Today's Schedule presents each appointment with contextual action buttons,
while a searchable Patient List Sidebar allows rapid access to any patient record. A
priority-based Lab Reports Queue ensures that critical results are reviewed first. The
Prescription Writing Interface incorporates drug interaction checking to support safe
prescribing, and a Leave Management Widget allows doctors to submit and track leave
applications without leaving the dashboard.

\subsection{Admin Dashboard Module}

\justifying
The admin dashboard provides hospital administrators with tools for managing facility
operations across all clinical and administrative functions. Statistics Cards give a
real-time overview of key performance indicators, and a Quick Action Cards Grid exposes
the most common administrative tasks in a single click. A Doctor Management Interface
supports inline editing of doctor profiles and schedules, while a Medicine Inventory
Management panel provides visibility into stock levels, surfaces low-stock alerts, and
tracks expiry dates to prevent medication shortages. A Lab Technician Management list
allows administrators to onboard and manage laboratory staff, and a Leave Approval
Workflow with calendar view enables efficient review and disposition of leave requests.

\subsection{Lab Technician Module}

\justifying
The laboratory technician module streamlines test processing from patient identification
through result entry to report submission. A multi-step Send Report Form guides the
technician through each stage of the reporting process, reducing the likelihood of
incomplete submissions. A searchable Patient History Table gives technicians immediate
access to prior results for the same patient, supporting trend analysis and anomaly
detection. A Statistics Dashboard displays daily report counts and turnaround times,
helping laboratory supervisors monitor throughput against service-level targets. A Quality
Control Panel tracks equipment status and reagent lot information, ensuring that results
are always produced under validated conditions.

\subsection{Appointment Booking Module}

\justifying
The appointment booking module provides an intuitive interface for finding and scheduling
appointments. A Doctor Search Interface with autocomplete and advanced filters allows
patients to narrow candidates by specialization, location, language, and availability. A
tabbed Doctor Profile View presents qualifications, reviews, and fee information in an
organized layout. An Availability Calendar uses color-coded slots to communicate at a
glance which times are open, and a real-time Time Slot Selection grid locks the selected
slot during the booking process to prevent concurrent bookings of the same time. A Booking
Form captures symptoms and insurance details before the patient proceeds to payment, and a
Confirmation Modal handles payment processing and offers calendar integration so the
appointment is automatically added to the patient's personal calendar.

% =============================================================================
% 4.6 SECURITY ARCHITECTURE
% =============================================================================
\section{Security Architecture}

\justifying
The BondHealth platform implements a defense-in-depth security architecture embedded at every
layer of the system. The guiding principles are Defense in Depth, Least Privilege, Secure by
Default, and Complete Mediation, which together ensure that every access attempt is
authenticated, authorised, and logged regardless of the path by which a resource is reached.
Rather than relying on a single perimeter control, security measures are applied at the
network boundary, the application layer, the database tier, and the file storage layer
simultaneously, so that the compromise of any one control does not expose the system as a
whole. The table below summarises the security measures implemented at each layer.

\begin{table}[H]
    \small
    \centering
    \renewcommand{\arraystretch}{1.3}
    \begin{tabular}{|>{\raggedright\arraybackslash}p{2.8cm}|>{\raggedright\arraybackslash}p{3.8cm}|>{\raggedright\arraybackslash}p{6.8cm}|}
        \hline
        \textbf{Security Layer} & \textbf{Implementation} & \textbf{Description} \\
        \hline
        Authentication
            & JWT + HTTP-only cookies
            & Tokens with HttpOnly, Secure flags. 7-day expiry. \\
        \hline
        Password Security
            & Bcrypt (10 rounds)
            & Salted hashes. Rate limiting: 5 attempts then lockout. \\
        \hline
        Database Security
            & SSL/TLS + Neon PostgreSQL
            & TLS 1.2+ encryption. Credentials in environment variables. \\
        \hline
        Session Management
            & Redis server-side store
            & Auto expiration. Max 5 concurrent sessions. \\
        \hline
        Role-Based Access
            & RBAC middleware
            & Permission inheritance. Hospital-scoped isolation. \\
        \hline
        Data Encryption
            & HTTPS/TLS 1.3
            & HSTS headers. Perfect forward secrecy. \\
        \hline
        Input Validation
            & JSON schema validation
            & SQL injection prevention via parameterized queries. \\
        \hline
        Audit Logging
            & DB triggers
            & 7-year retention. Logs auth attempts and data access. \\
        \hline
        File Security
            & ClamAV scanning
            & Type validation by magic numbers. Size limits enforced. \\
        \hline
    \end{tabular}
    \caption{Security Features Summary}
    \label{tab:ch4-security}
\end{table}

\subsection{Data Encryption}

\justifying
Data at rest is protected at multiple levels. At the storage layer, Transparent Data
Encryption (TDE) with AES-256 ensures that database files on disk are unreadable without
the appropriate key. Sensitive fields containing personally identifiable information and
medical data receive an additional layer of application-level encryption using AES-256-GCM
before they are written to the database, so that even a privileged database user cannot read
the raw values. Files uploaded to object storage are encrypted server-side using AWS KMS or
Azure Key Vault, depending on the deployment environment.

\noindent
\justifying
Data in transit is equally protected. All external communications require TLS 1.3 as a
minimum, and database connections enforce TLS with full certificate validation to guard
against man-in-the-middle attacks. Every API endpoint is served with HTTP Strict Transport
Security preload headers and is configured for perfect forward secrecy, ensuring that the
compromise of a long-term private key cannot be used to decrypt previously recorded
sessions.

\subsection{Compliance and Audit}

\justifying
The platform is designed to satisfy HIPAA requirements for the protection of protected health
information (PHI). Access controls and audit trails cover all PHI, and all such data is
encrypted both at rest and in transit as described above. In the event of a breach, procedures
are in place to notify affected parties within 72 hours. Sessions are automatically terminated
after 15 minutes of inactivity to reduce the risk of unauthorized access on unattended devices.

\noindent
\justifying
All security-relevant events are written to an immutable audit log. Authentication attempts,
including failures, are recorded with the originating IP address and user agent. Every access
to a PHI record is logged with the patient identifier and the type of record accessed. Data
modifications capture both the before and after values to support forensic reconstruction of
any change. Log records are retained for seven years to satisfy regulatory requirements across
all applicable jurisdictions.

% =============================================================================
% 4.7 DATABASE CONNECTION DESIGN
% =============================================================================

\section{Database Connection Design (Neon PostgreSQL)}

\justifying
The database connection architecture is designed for optimal performance, reliability,
and security when connecting the Node.js application layer to Neon PostgreSQL. Because each
incoming HTTP request may require several database queries, managing connections efficiently
is critical to sustaining acceptable response times under load. The architecture addresses
this through connection pooling, strict SSL configuration, and disciplined credential
management.

\subsection{Connection Pooling Architecture}

\justifying
Connection pooling maintains a cache of reusable database connections, eliminating the
overhead of a full TCP handshake per request. The pool maintains a configurable minimum of
ready connections to handle sudden traffic spikes gracefully, and automatically detects and
replaces failed or stale connections without manual intervention. The pool is configured with
a minimum size of 5 connections to absorb sudden spikes, a maximum size of 20 connections to
limit load on the database server, a connection timeout of 10 seconds, an idle timeout of
30 seconds after which unused connections are released, and a maximum connection lifetime of
30 minutes to ensure stale connections are periodically recycled. Together these settings
reduce latency, improve resource utilization, provide automatic recovery from transient
failures, and smooth out the effects of bursty traffic patterns.

\begin{lstlisting}[language=JavaScript,
    caption={Database Connection Configuration (config/database.js)},
    label={lst:dbconfig},
    breaklines=true,
    columns=flexible,
    keepspaces=true]
const { Pool } = require('@neondatabase/serverless');

const pool = new Pool({
    connectionString: process.env.DATABASE_URL,
    ssl: { require: true, rejectUnauthorized: true },
    max: 20,
    idleTimeoutMillis: 30000,
    connectionTimeoutMillis: 10000,
});

const query = async (text, params) => {
    const start = Date.now();
    try {
        const res = await pool.query(text, params);
        const duration = Date.now() - start;
        if (duration > 100) {
            console.warn('Slow query detected:', { text, duration });
        }
        return res;
    } catch (error) {
        console.error('Database query error:', error);
        throw error;
    }
};
\end{lstlisting}

\subsection{Connection Security}

\justifying
All connections to Neon PostgreSQL are established with SSL mode set to \texttt{require},
which prevents any unencrypted connection from succeeding. Certificate validation is enabled
on the client side to protect against man-in-the-middle attacks, and the minimum acceptable
TLS version is TLSv1.2. The database connection string is stored exclusively in environment
variables and is never committed to version control. Credentials are rotated on a quarterly
schedule and immediately after any security incident, limiting the window of exposure in the
event that a credential is compromised.

\justifying
\noindent
Beyond transport-level protection, the application connects to the database using a
dedicated service account granted only the privileges required by the application's query
set — namely \texttt{SELECT}, \texttt{INSERT}, \texttt{UPDATE}, and \texttt{DELETE} on
the relevant tables, with no \texttt{DDL} privileges. This ensures that even if the
application-layer credentials were obtained by an attacker, they could not be used to
alter the schema, drop tables, or escalate access within the database. Query parameters
are always passed separately from the SQL statement text, so the database driver handles
escaping independently of application code, eliminating an entire class of injection
vulnerability at the connection level rather than relying solely on application-layer
validation.

% =============================================================================
% 4.8 THIRD-PARTY INTEGRATIONS
% =============================================================================

\section{Third-party Integrations}

\justifying
The BondHealth platform integrates with multiple external services to extend functionality
while keeping development focused on core healthcare features. All integrations communicate
over HTTPS with credentials stored exclusively in environment variables, and each integration
is abstracted behind a thin service wrapper so that the underlying provider can be swapped
without changes to the calling code. Where an external service supports webhook delivery,
the platform registers a dedicated endpoint to receive asynchronous event notifications,
decoupling the user-facing response from the outcome of the background operation. Each
integration is also subject to the same error handling and retry conventions as internal
modules, so that a transient failure in a third-party service degrades gracefully rather
than propagating an unhandled exception to the user. The table below summarizes each
integration, its purpose, and key implementation details.

\begin{table}[H]
    \small
    \centering
    \renewcommand{\arraystretch}{1.3}
    \begin{tabular}{|>{\raggedright\arraybackslash}p{3.0cm}|>{\raggedright\arraybackslash}p{2.8cm}|>{\raggedright\arraybackslash}p{7.6cm}|}
        \hline
        \textbf{Integration} & \textbf{Purpose} & \textbf{Implementation} \\
        \hline
        Stripe/\allowbreak Razorpay
            & Payment gateway
            & Webhooks for async notifications. PCI DSS Level 1 compliant. \\
        \hline
        Twilio/\allowbreak MessageBird
            & SMS alerts
            & Templates for reminders. Delivery tracking enabled. \\
        \hline
        SendGrid/\allowbreak NodeMailer
            & Email service
            & Transactional emails. DKIM/SPF configured. \\
        \hline
        AWS S3/\allowbreak Azure Blob
            & File storage
            & AES-256 encryption. Signed URLs (5-min expiry). \\
        \hline
        Socket.io
            & Real-time chat
            & Redis adapter for scaling. Room-based messaging. \\
        \hline
        Google Calendar API
            & Calendar sync
            & OAuth 2.0. Create/update events with reminders. \\
        \hline
        Firebase Cloud Messaging
            & Push notifications
            & Cross-platform support. Topic-based subscriptions. \\
        \hline
        Google Maps API
            & Location services
            & Distance calculation. Address autocomplete. \\
        \hline
    \end{tabular}
    \caption{Third-party Integrations}
    \label{tab:integrations}
\end{table}

% =============================================================================
% 4.9 PERFORMANCE OPTIMIZATION
% =============================================================================

\section{Performance Optimization}

\justifying
The BondHealth platform implements multiple optimization strategies across all layers of the
stack to ensure responsive user experiences under high load. These strategies address four
complementary areas: database query performance, in-memory caching, API and frontend delivery,
and infrastructure-level scaling and monitoring.

\justifying
\noindent
At the database level, strategic indexing on frequently queried columns reduces query execution
times significantly, and connection pooling eliminates per-request connection overhead as
described in Section~4.7. Query plans are analyzed using \texttt{EXPLAIN} to identify missing
indexes or suboptimal join strategies, and reporting queries are routed to read replicas to
offload the primary database. Large tables such as appointments and logs are partitioned by date
and month respectively, keeping individual partition sizes manageable as data accumulates over
time.

\justifying
\noindent
Caching is implemented using Redis at multiple levels. Frequently read but slowly changing
data --- such as doctor lists and hospital information --- is cached in Redis with a
configurable TTL, while session data is stored in Redis for sub-millisecond access on every
authenticated request. Public endpoints benefit from response caching served through a CDN,
and the results of expensive but stable database queries are cached to prevent redundant
computation on repeated calls.

\justifying
\noindent
At the API and frontend layer, all list endpoints use pagination to keep response payloads
within predictable bounds. Gzip and Brotli compression reduce response sizes by 70 to 80
percent, and HTTP/2 multiplexing allows the browser to issue parallel requests over a single
connection. A batch endpoint is provided for mobile clients that need to retrieve multiple
resources in a single round trip. On the frontend, images are lazy-loaded, JavaScript bundles
are code-split to avoid shipping unused code, and static assets are served from a CDN. A
service worker enables offline capability and client-side caching of recently viewed data.

\justifying
\noindent
At the infrastructure level, the application is designed to be stateless, allowing horizontal
scaling behind an NGINX load balancer without session affinity requirements. Auto-scaling groups
respond to CPU and memory metrics to add or remove instances as demand fluctuates. A CDN in
front of static assets reduces origin load by approximately 90 percent. Real-time performance
monitoring via DataDog or New Relic triggers alerts when the 95th-percentile latency exceeds
500 milliseconds, and synthetic monitoring continuously exercises critical user journeys to
detect regressions before real users are affected. Log aggregation pipelines surface slow query
patterns for ongoing tuning.

% =============================================================================
% 4.10 SUMMARY
% =============================================================================

\section{Summary}

This chapter has presented the complete system design of the BondHealth platform,
spanning the three-tier architecture, data flow diagrams, database schema, API design,
module specifications, security architecture, and performance optimisation strategies.
The design decisions reflect consistent priorities: security of sensitive patient data,
maintainability of a multi-role codebase, and scalability for growing institutional
adoption. The layered separation between presentation, business logic, and data persistence
ensures each tier can evolve independently. The database schema, anchored by UUID primary
keys and enforced through foreign key constraints, provides a reliable foundation for the
platform's clinical workflows. The RESTful API exposes this functionality in a predictable,
role-scoped manner, while the security architecture ensures every access attempt is
authenticated, authorised, and logged. The following chapter details the implementation of
each module described here, translating these design specifications into working code.